% Omslagmotief scheikunde: structuurformules (zeshoeken) en een kolf.
\foreach \cx/\cy/\k in {10/38/kleur1,30.78/38/kleur2,20.39/20/kleur3,51.56/38/kleur1} {
  \draw[\k!60!white,line width=1.5pt] (\cx,\cy) +(30:12) -- +(90:12) -- +(150:12) -- +(210:12) -- +(270:12) -- +(330:12) -- cycle;}
\draw[kleur1!60!white,line width=1.2pt] (10,38) circle (7);
\draw[white,opacity=.6,line width=1pt] (61.952,44) -- (68.881,48) node[right,font=\sffamily\small,opacity=.9] {OH};
\draw[white,opacity=.6,line width=1pt] (20.39,8) -- (20.39,0) node[below,font=\sffamily\small,opacity=.9] {NH$_2$};
\begin{scope}[shift={(96,4)}]
  \draw[kleur4!65!white,line width=1.5pt] (-4,40) -- (-4,24) -- (-16,0) arc (200:340:17 and 6) -- (4,24) -- (4,40);
  \fill[kleur2!40!white,opacity=.35] (-12,8) -- (-16,0) arc (200:340:17 and 6) -- (12,8) -- cycle;
  \fill[white] (-3,12) circle (1); \fill[white] (2,16) circle (.8); \fill[white] (-1,20) circle (.6);
\end{scope}
